\documentclass[11pt]{article}

\usepackage[T1]{fontenc}
\usepackage[utf8]{inputenc}
\usepackage[margin=1in]{geometry}
\usepackage{amsmath,amssymb,amsthm}
\usepackage{booktabs}
\usepackage{array}
\usepackage[hidelinks]{hyperref}

\setlength{\emergencystretch}{1.5em}

\newtheorem{theorem}{Theorem}[section]
\newtheorem{proposition}[theorem]{Proposition}
\newtheorem{lemma}[theorem]{Lemma}
\newtheorem{corollary}[theorem]{Corollary}
\theoremstyle{definition}
\newtheorem{definition}[theorem]{Definition}
\newtheorem{remark}[theorem]{Remark}

\newcommand{\Cl}{\mathrm{Cl}}
\newcommand{\End}{\mathrm{End}}
\newcommand{\rad}{\mathrm{rad}}
\newcommand{\Nbar}{N_{\mathrm{bar}}}
\newcommand{\Nrev}{N_{\mathrm{rev}}}
\newcommand{\Nreg}{N_{\mathrm{reg}}}

\title{\textbf{Disproof of Conjecture 00000007622:\\
Zero Divisors, Involution Norms, Isotropic Vectors,\\
and Minimal Left Ideals in Clifford Algebras}}

\author{Formal Verification \& Mathematical Analysis Report}
\date{November 2026}

\begin{document}
\maketitle

\begin{abstract}
We resolve Conjecture~\texttt{00000007622} in the negative across every standard
Clifford algebra convention. The conjecture asserts three linked claims for a
Clifford algebra $\Cl(V, Q)$ associated to a quadratic space $(V, Q)$:
(A)~vanishing of the conversion/Clifford norm is necessary and sufficient for a
nonzero multivector to be a zero divisor; (B)~the kernel of necessity is given by
isotropic vectors of $(V, Q)$; and (C)~the minimal left ideal of zero divisors is
the spinor module. We first isolate the three restricted regimes that motivate the
conjecture---grade-$1$ vectors in arbitrary dimension, full multivectors in
dimension $\dim V \le 2$, and primitive-idempotent ideals in simple split algebras
such as $\Cl_{1,1}(F) \cong M_2(F)$---and then prove unconditionally that all three
clauses fail in general. Specifically: (1)~strictly anisotropic Euclidean spaces
yield zero divisors under both defining relations $v^2 = +Q(v)\cdot 1$ (already in
$\Cl_{1,0}(F) \cong F \oplus F$, even over $F = \mathbb{C}$) and $v^2 = -Q(v)\cdot 1$
(in $\Cl_{0,3}(\mathbb{R}) \cong \mathbb{H} \oplus \mathbb{H}$), while even in
$\Cl_{1,1}(F)$ the annihilator kernel of $1 + e_0$ meets $V$ trivially;
(2)~in dimension $3$ ($\Cl_{0,3}, \Cl_{2,1}$) and universally in dimension $4$
(both $\Cl_{4,0}$ and $\Cl_{0,4}$), $x = 1 + \omega$ is a two-sided zero divisor
fixed by all canonical involutions, so $x\,\iota(x) = 2x \neq 0$ and
$\langle x\,\iota(x)\rangle_0 = 2 \neq 0$; and (3)~minimal left ideals fail to be
spinor modules in degenerate exterior algebras $\Lambda(F^n)$, split semisimple
Clifford algebras, and division Clifford algebras. All results are verified by
exact rational computation (\texttt{verify.py}) and machine-checked in Core
Lean~4.19.0 with zero external dependencies and zero \texttt{sorry} axioms.
\end{abstract}

\section{Introduction and Statement of the Conjecture}

Conjecture~\texttt{00000007622} from \emph{The Last Math Competition} is stated
bilingually as follows:
\begin{quote}
\textbf{English.} \emph{Definition: Zero divisors with the divisor structure of
Clifford algebras. Conjecture: The conversion constant is the zero-divisor
criterion and the norm-vanishing criterion is necessary and sufficient; the kernel
of necessity is the isotropic vector of the quadratic form; and the minimal left
ideal of zero divisors is the spinor module. (norm-vanishing conversion)}
\end{quote}
In the accompanying Chinese formulation, the three clauses read:
(i)~the conversion constant serves as the zero-divisor criterion and the
norm-vanishing criterion is necessary and sufficient;
(ii)~the kernel of the necessity direction is the isotropic vector of the
quadratic form; and
(iii)~the minimal left ideal of zero divisors is the spinor module.

To ensure that our negative resolution does not depend on any ambiguity of
terminology or sign convention, we analyze every standard algebraic formulation
of Clifford algebras over a field $F$ with $\mathrm{char}(F) \neq 2$ (including
$F = \mathbb{Q}$, $F = \mathbb{R}$, and $F = \mathbb{C}$):
\begin{enumerate}
  \item \textbf{Both defining sign conventions:} the algebraic/Chevalley convention
    $\Cl^+(V, Q)$ defined by $v^2 = +Q(v)\cdot 1$, and the geometric/Lawson--Michelsohn
    convention $\Cl^-(V, Q)$ defined by $v^2 = -Q(v)\cdot 1$.
  \item \textbf{All canonical involutions and norms:} the multivector Clifford
    conjugation norm $\Nbar(x) = x\bar{x}$, the reversion norm $\Nrev(x) = x\tilde{x}$,
    the grade-involution product $x\,\alpha(x)$, their grade-$0$ scalar trace
    projections $\langle x\,\iota(x)\rangle_0 \in F$, and the regular-representation
    determinant $\Nreg(x) = \det(L_x) \in F$.
  \item \textbf{Both readings of isotropic necessity:} existence of nonzero
    isotropic vectors in $(V, Q)$ as a prerequisite for zero divisors in
    $\Cl(V, Q)$, as well as element-level intersection $\ker(L_x) \cap V$ for
    individual zero divisors $x$.
  \item \textbf{All regimes of $(V, Q)$:} positive-definite, indefinite/split,
    and degenerate (including totally degenerate exterior algebras $\Lambda(V)$).
\end{enumerate}

\textbf{Main Verdict.} Conjecture~\texttt{00000007622} is \textbf{FALSE} under
every combination of the conventions above.

\section{Algebraic Preliminaries and Graded Construction}

Let $F$ be a field with $\mathrm{char}(F) \neq 2$, and let $V$ be an
$n$-dimensional $F$-vector space equipped with a quadratic form $Q : V \to F$ and
associated symmetric bilinear form $B(u, v) = \frac{1}{2}\bigl(Q(u + v) - Q(u) - Q(v)\bigr)$.
Every quadratic form over $F$ ($\mathrm{char}(F) \neq 2$) admits a $B$-orthogonal
basis $e_0, e_1, \dots, e_{n-1}$ with $Q(e_i) = c_i \in F$, so that
$Q\bigl(\sum_{i=0}^{n-1} v_i e_i\bigr) = \sum_{i=0}^{n-1} c_i v_i^2$.

\begin{definition}[Clifford algebras $\Cl^+(V, Q)$ and $\Cl^-(V, Q)$]
For $\sigma = (c_0, \dots, c_{n-1}) \in F^n$, the Clifford algebra
$\Cl(F^n, \sigma) := T(V) / \langle v \otimes v - Q(v)\cdot 1 : v \in V\rangle$
has dimension $2^n$ with ordered blade basis
$e_S = e_{i_1} e_{i_2} \cdots e_{i_k}$ indexed by subsets
$S = \{i_1 < \dots < i_k\} \subseteq \{0, \dots, n-1\}$, governed by
$e_i^2 = c_i \cdot 1$ and $e_i e_j = -e_j e_i$ ($i \neq j$). Writing
$-\sigma = (-c_0, \dots, -c_{n-1})$:
\[
  \Cl^+(V, Q) := \Cl(F^n, \sigma) \quad (v^2 = +Q(v)\cdot 1), \qquad
  \Cl^-(V, Q) := \Cl(F^n, -\sigma) \quad (v^2 = -Q(v)\cdot 1).
\]
Over $F$ with Euclidean form $Q_{\mathrm{pos}}(v) = \sum_{i=0}^{n-1} v_i^2$,
$\Cl^+(F^n, Q_{\mathrm{pos}}) = \Cl_{n,0}(F)$ has $e_i^2 = +1$, while
$\Cl^-(F^n, Q_{\mathrm{pos}}) = \Cl_{0,n}(F)$ has $e_i^2 = -1$.
\end{definition}

\begin{proposition}[Recursive $\mathbb{Z}/2\mathbb{Z}$-graded model]
\label{prop:recursive}
For any commutative unital ring $R$ and $\sigma = (c_0, c_1, \dots, c_{n-1}) \in R^n$,
let $A = \Cl(R^{n-1}, (c_1, \dots, c_{n-1}))$ with grade involution $\alpha : A \to A$.
Then $\Cl(R^n, \sigma) \cong A \oplus e_0 A$ via $(a, b) \leftrightarrow a + e_0 b$,
with associative unital multiplication and involutions given recursively by
\begin{align}
  (a + e_0 b)(u + e_0 v) &= \bigl(a u + c_0\,\alpha(b) v\bigr) + e_0\bigl(\alpha(a) v + b u\bigr), \\
  \alpha(a + e_0 b) &= \alpha(a) - e_0\,\alpha(b), \qquad
  \widetilde{a + e_0 b} = \tilde{a} + e_0\,\alpha(\tilde{b}), \qquad
  \overline{a + e_0 b} = \alpha(\widetilde{a + e_0 b}).
\end{align}
\end{proposition}
\begin{proof}
Since $e_0$ anticommutes with $e_1, \dots, e_{n-1}$, moving any $w \in A$ past
$e_0$ yields $w e_0 = e_0\,\alpha(w)$, and $e_0^2 = c_0$. Expanding
$(a + e_0 b)(u + e_0 v) = au + a e_0 v + e_0 b u + e_0 b e_0 v$ substitutes
$a e_0 = e_0\alpha(a)$ and $e_0 b e_0 = e_0^2 \alpha(b) = c_0 \alpha(b)$.
Associativity follows directly from $\alpha^2 = \mathrm{id}_A$ and the
automorphism property $\alpha(xy) = \alpha(x)\alpha(y)$: comparing the even and
odd components of $((a,b)(u,v))(p,q)$ and $(a,b)((u,v)(p,q))$ gives identical
four-term sums. Similarly, $(a + e_0 b)^\sim = \tilde{a} + \tilde{b} e_0 =
\tilde{a} + e_0\alpha(\tilde{b})$ reverses products by induction.
\end{proof}

On an orthogonal $k$-blade $e_S$ ($|S| = k$), the three canonical involutions act
by pure grade signs:
\begin{equation}
\label{eq:grade-signs}
  \alpha(e_S) = (-1)^k e_S, \qquad
  \tilde{e}_S = (-1)^{k(k-1)/2} e_S, \qquad
  \bar{e}_S = (-1)^{k(k+1)/2} e_S.
\end{equation}
In particular, Clifford conjugation $\bar{x}$ negates grades $k \equiv 1, 2 \pmod 4$
and \emph{fixes} grades $k \equiv 0, 3 \pmod 4$, whereas reversion $\tilde{x}$
negates grades $k \equiv 2, 3 \pmod 4$ and \emph{fixes} grades $k \equiv 0, 1 \pmod 4$.

\begin{definition}[Zero divisors, isotropic vectors, and minimal left ideals]
Let $A = \Cl(V, Q)$ over a field $F$.
\begin{enumerate}
  \item An element $x \in A \setminus \{0\}$ is a \emph{left zero divisor} if
    $\exists y \in A \setminus \{0\}$ with $xy = 0$, and a \emph{right zero divisor}
    if $\exists y \in A \setminus \{0\}$ with $yx = 0$. Since $\dim_F A = 2^n < \infty$,
    left multiplication $L_x : A \to A$ is injective iff surjective iff invertible:
    \begin{align*}
      x \neq 0 \text{ is a left zero divisor}
      &\iff x \neq 0 \text{ is a right zero divisor} \\
      &\iff x \notin A^\times \iff \det(L_x) = 0.
    \end{align*}
  \item A vector $v \in V$ is \emph{isotropic} if $v \neq 0$ and $Q(v) = 0$. The
    quadratic space $(V, Q)$ is \emph{anisotropic} if no isotropic vector exists.
  \item A nonzero left ideal $I \subseteq A$ is \emph{minimal} if every nonzero
    left ideal $J \subseteq I$ satisfies $J = I$. An \emph{algebraic spinor module}
    of a simple nondegenerate Clifford algebra $A \cong M_r(D)$ ($D$ a division
    $F$-algebra) is a minimal left ideal $S = A f$ generated by a primitive
    idempotent $f^2 = f \neq 0$; such an $S$ is a simple, faithful left $A$-module.
\end{enumerate}
\end{definition}

\section{Restricted Positive Regimes Motivating the Conjecture}

Before refuting Conjecture~\texttt{00000007622} in full generality, we establish
three exact positive theorems showing where each clause holds under strict
dimensional or grade hypotheses. These theorems also serve as constructive
non-vacuity witnesses in our Lean~4 formalization.

\begin{theorem}[Grade-$1$ vector criterion in arbitrary dimension]
\label{thm:grade1}
Let $(V, Q)$ be any finite-dimensional quadratic space over a field $F$, and let
$v \in V \subset \Cl^\pm(V, Q)$. Then $v$ is a zero divisor in $\Cl^\pm(V, Q)$
if and only if $v$ is an isotropic vector of $(V, Q)$ ($v \neq 0$ and $Q(v) = 0$).
\end{theorem}
\begin{proof}
Under $v^2 = \pm Q(v)\cdot 1$, if $v \neq 0$ and $Q(v) = 0$, then $v \cdot v = 0$,
witnessing that $v$ is a two-sided zero divisor. Conversely, if $v = 0$, $v$ is
not a zero divisor by definition; if $Q(v) \neq 0$, then $u := (\pm Q(v))^{-1} v$
satisfies $vu = uv = 1$, so $v \in \Cl^\pm(V, Q)^\times$ cannot be a zero divisor.
\end{proof}

\begin{theorem}[Low-dimensional Clifford norm criterion, $\dim V \le 2$]
\label{thm:dim2}
Let $(V, Q)$ have orthogonal signature $\sigma$ of dimension $n \le 2$ over a
field $F$. Then:
\begin{enumerate}
  \item For every $x \in \Cl(F^n, \sigma)$, the Clifford conjugation product is
    scalar: $x\bar{x} = \bar{x}x = \Nbar(x)\cdot 1$ with $\Nbar(x) \in F$. Explicitly,
    for $n = 2$ and $x = a_0 + a_1 e_0 + a_2 e_1 + a_{12} e_0 e_1$,
    \[
      \Nbar(x) = a_0^2 - c_0 a_1^2 - c_1 a_2^2 + c_0 c_1 a_{12}^2, \qquad
      \det(L_x) = \Nbar(x)^2.
    \]
  \item Consequently, for every $x \in \Cl(F^n, \sigma) \setminus \{0\}$,
    $x$ is a zero divisor if and only if $x\bar{x} = 0$ (equivalently $\Nbar(x) = 0$).
  \item Over any ordered field $F$ (in particular $\mathbb{Q}$ and $\mathbb{R}$),
    $\Cl_{0,1}(F) = \Cl(F, [-1])$ and $\Cl_{0,2}(F) = \Cl(F, [-1, -1])$ are
    division algebras with zero zero-divisors.
\end{enumerate}
\end{theorem}
\begin{proof}
For $n \le 2$, only grades $0, 1, 2$ occur. By~\eqref{eq:grade-signs},
$\bar{e}_0 = -e_0$, $\bar{e}_1 = -e_1$, and $\overline{e_0 e_1} = -e_0 e_1$.
Writing $x = a_0 + m$ with $m = a_1 e_0 + a_2 e_1 + a_{12} e_0 e_1$, we have
$\bar{x} = a_0 - m$, and pairwise anticommutation among $\{e_0, e_1, e_0 e_1\}$
together with $(e_0 e_1)^2 = -c_0 c_1$ yields
$x\bar{x} = \bar{x}x = a_0^2 - m^2 = (a_0^2 - c_0 a_1^2 - c_1 a_2^2 + c_0 c_1 a_{12}^2)\cdot 1$.
For $x \neq 0$, involution bijectivity ($\bar{\bar{x}} = x$) ensures $\bar{x} \neq 0$.
Thus if $\Nbar(x) = 0$, $x\bar{x} = 0$ directly exhibits $x$ as a zero divisor;
if $\Nbar(x) \neq 0$, then $\Nbar(x)^{-1}\bar{x}$ is a two-sided inverse of $x$.
When $c_0 = c_1 = -1$ over an ordered field, $\Nbar(x) = a_0^2 + a_1^2 + a_2^2 + a_{12}^2$
is a sum of squares, vanishing only at $x = 0$.
\end{proof}

\begin{theorem}[Simple faithful spinor ideal in $\Cl_{1,1}(F) \cong M_2(F)$]
\label{thm:cl11}
Over any field $F$ with $\mathrm{char}(F) \neq 2$, let $\Cl_{1,1}(F) = \Cl(F^2, [1, -1])$
with $e_0^2 = 1$, $e_1^2 = -1$, and primitive idempotent $f = \frac{1}{2}(1 + e_0)$.
Then $S = \Cl_{1,1}(F) f$ is a $2$-dimensional minimal left ideal, every
$x \in S \setminus \{0\}$ is a zero divisor, and $S$ is a faithful simple left
$\Cl_{1,1}(F)$-module realizing $\Cl_{1,1}(F) \xrightarrow{\sim} \End_F(S) \cong M_2(F)$.
\end{theorem}
\begin{proof}
Since $e_0^2 = 1$, $f^2 = \frac{1}{4}(2 + 2e_0) = f$, and $f(1 - f) = 0$ with
$1 - f = \frac{1}{2}(1 - e_0) \neq 0$, so every $x = af \in S \setminus \{0\}$
satisfies $x(1 - f) = 0$. Writing $a = u + e_0 v$ with $u, v \in \Cl(F, [-1])$ gives
$af = w + e_0 \alpha(w)$ where $w = \frac{1}{2}(u + \alpha(v)) = p + q e_1$, so
$\dim_F S = 2$. Given any nonzero $x = w + e_0 \alpha(w)$ in a nonzero left
subideal $J \subseteq S$: if $p \neq 0$, then $x + e_0 x = 4p f \in J$ implies
$f \in J$; if $p = 0$, then $q \neq 0$ and $e_1 x = -q + e_0(-q) = -2q f \in J$
again implies $f \in J$, so $J = S$ and $S$ is minimal. Finally, if
$(u + e_0 v) S = 0$, evaluating on $2f = 1 + e_0$ and $2 e_1 f = e_1 - e_0 e_1$
yields $u + \alpha(v) = 0$ and $(u - \alpha(v))e_1 = 0$; multiplying the latter
by $-e_1$ ($e_1(-e_1) = 1$) gives $u - \alpha(v) = 0$, whence $2u = 0 \implies u = v = 0$.
\end{proof}

\section{Disproof of Clause B: Zero Divisors Without Isotropic Vectors}

Clause~B fails catastrophically as soon as one passes from grade-$1$ vectors
(Theorem~\ref{thm:grade1}) to mixed-grade multivectors, under \emph{both} sign
conventions and over $\mathbb{Q}$, $\mathbb{R}$, and $\mathbb{C}$.

\begin{theorem}[Anisotropic zero divisors under both sign conventions]
\label{thm:disproof-B}
\leavevmode
\begin{enumerate}
  \item \textbf{Convention $v^2 = +Q(v)\cdot 1$ (arbitrary field $F$, $\mathrm{char}(F) \neq 2$):}
    On the $1$-dimensional quadratic space $(F^1, Q(v_0) = v_0^2)$, $Q(v_0) = 0 \iff v_0 = 0$,
    so $(F^1, Q)$ has \emph{zero} nonzero isotropic vectors (even when $F = \mathbb{C}$).
    Yet in $\Cl^+(F^1, Q) = \Cl_{1,0}(F) \cong F \oplus F$, the nonzero elements
    $x = 1 + e_0$ and $y = 1 - e_0$ satisfy $xy = yx = 1 - e_0^2 = 0$.
    Over any ordered field $F$ ($\mathbb{Q}, \mathbb{R}$), the same holds for all
    Euclidean spaces $(F^n, \sum_{i=0}^{n-1} v_i^2)$ with $n \ge 1$ (including
    simple $\Cl_{2,0}(\mathbb{R}) \cong M_2(\mathbb{R})$).
  \item \textbf{Convention $v^2 = -Q(v)\cdot 1$ (ordered fields $\mathbb{Q}, \mathbb{R}$, and $\mathbb{C}$):}
    Over any ordered field $F$, the $3$-dimensional Euclidean space
    $(F^3, v_0^2 + v_1^2 + v_2^2)$ is strictly anisotropic ($Q(v) = 0 \implies v = 0$).
    Yet in $\Cl^-(F^3, Q_{\mathrm{pos}}) = \Cl_{0,3}(F)$ ($e_i^2 = -1$), the central
    volume blade $\omega_3 = e_0 e_1 e_2$ satisfies
    \[
      \omega_3^2 = (e_0 e_1 e_2)(e_0 e_1 e_2) = (-1)^3 e_0^2 e_1^2 e_2^2
                 = (-1)(-1)^3 = +1,
    \]
    so $(1 + \omega_3)(1 - \omega_3) = 1 - \omega_3^2 = 0$ with $1 \pm \omega_3 \neq 0$.
    (Over $F = \mathbb{C}$, even $n = 1$ under $v^2 = -Q(v)$ has $(i + e_0)(i - e_0) = -1 - (-1) = 0$.)
  \item \textbf{Element-level annihilator kernel in $\Cl_{1,1}(F)$:}
    Even in split signature $\Cl_{1,1}(F)$ ($e_0^2 = 1, e_1^2 = -1$), where nonzero
    isotropic vectors $e_0 \pm e_1 \in V \setminus \{0\}$ exist, the zero divisor
    $x = 1 + e_0$ satisfies $\ker(L_{1 + e_0}) \cap V = \{0\}$.
\end{enumerate}
\end{theorem}
\begin{proof}
Parts~(1) and~(2) are immediate from the displayed products; strict anisotropy of
$Q_{\mathrm{pos}}$ over an ordered field follows from positivity of nonzero squares.
For part~(3), let $v = v_0 e_0 + v_1 e_1 \in V$. Then
$(1 + e_0)(v_0 e_0 + v_1 e_1) = v_0 + v_0 e_0 + v_1 e_1 + v_1 e_0 e_1$, whose
grade-$0$ and grade-$1$ coordinates vanish iff $v_0 = v_1 = 0$.
\end{proof}

\section{Disproof of Clause A: Breakdown of Clifford and Involution Norms}

Theorem~\ref{thm:dim2} showed that $\Nbar(x) = x\bar{x}$ characterizes zero
divisors for $\dim V \le 2$ because Clifford conjugation negates grades $1$ and $2$.
At grade $3$ and grade $4$, Equation~\eqref{eq:grade-signs} gives
$(-1)^{3\cdot 4/2} = +1$ and $(-1)^{4\cdot 5/2} = +1$, causing Clause~A to collapse.

\begin{theorem}[Three-dimensional failure of Clifford conjugation norm]
\label{thm:dim3-norm}
Over any ordered field $F$ ($\mathbb{Q}, \mathbb{R}$), let
$\Cl_{0,3}(F) = \Cl^-(F^3, Q_{\mathrm{pos}})$ with $\omega_3 = e_0 e_1 e_2$, and
set $x = 1 + \omega_3$. Then:
\begin{enumerate}
  \item $x$ is a left and right zero divisor: $x(1 - \omega_3) = (1 - \omega_3)x = 0$,
    with $\det(L_x) = 0$ and $\mathrm{rank}(L_x) = 4 < 8$;
  \item $\bar{x} = 1 + \bar{\omega}_3 = 1 + \omega_3 = x$, so the multivector
    Clifford norm is
    \[
      x\bar{x} = (1 + \omega_3)^2 = 2(1 + \omega_3) \neq 0,
    \]
    which is neither scalar ($x\bar{x} \notin F\cdot 1$) nor zero;
  \item the grade-$0$ scalar trace norm is $\langle x\bar{x}\rangle_0 = 2 \neq 0$.
\end{enumerate}
The identical failure occurs in indefinite signature $\Cl_{2,1}(F)$ ($e_0^2 = e_1^2 = 1$,
$e_2^2 = -1$), where again $\omega_3^2 = (-1)^3(1)(1)(-1) = +1$ and $\bar{\omega}_3 = \omega_3$.
\end{theorem}
\begin{proof}
By~\eqref{eq:grade-signs} at $k = 3$, $\bar{\omega}_3 = (-1)^6 \omega_3 = \omega_3$.
Since $\omega_3^2 = 1$, $(1 + \omega_3)^2 = 1 + 2\omega_3 + 1 = 2 + 2\omega_3$.
Because $\mathrm{char}(F) \neq 2$, both the grade-$0$ coefficient $2$ and the
grade-$3$ coefficient $2$ are nonzero.
\end{proof}

Even more strikingly, in dimension $4$ we obtain a \emph{universal} counterexample
that simultaneously defeats \emph{both} sign conventions $v^2 = +Q(v)$ and
$v^2 = -Q(v)$ on positive-definite Euclidean $4$-space and \emph{all three}
canonical involutions at once.

\begin{theorem}[Universal $4$-dimensional counterexample across all conventions]
\label{thm:dim4-universal}
Let $F$ be any field with $\mathrm{char}(F) \neq 2$, let $\epsilon \in \{+1, -1\}$,
and consider $\Cl(F^4, [\epsilon, \epsilon, \epsilon, \epsilon])$---which equals
$\Cl_{4,0}(F) = \Cl^+(F^4, Q_{\mathrm{pos}})$ for $\epsilon = +1$ and
$\Cl_{0,4}(F) = \Cl^-(F^4, Q_{\mathrm{pos}})$ for $\epsilon = -1$. Let
$\omega_4 = e_0 e_1 e_2 e_3$ and $x = 1 + \omega_4$. Then:
\begin{enumerate}
  \item $\omega_4^2 = (-1)^{4\cdot 3/2} \epsilon^4 = (+1)(+1) = +1$, so
    $x(1 - \omega_4) = (1 - \omega_4)x = 0$ with $1 \pm \omega_4 \neq 0$, making $x$
    a two-sided zero divisor ($\det(L_x) = 0$, $\mathrm{rank}(L_x) = 8 < 16$);
  \item for \textbf{every} canonical involution $\iota \in \{\bar{\cdot}, \tilde{\cdot}, \alpha\}$,
    \[
      \iota(x) = x, \qquad
      x\,\iota(x) = 2(1 + \omega_4) \neq 0, \qquad
      \langle x\,\iota(x)\rangle_0 = 2 \neq 0;
    \]
  \item when $F$ is an ordered field ($\mathbb{Q}$ or $\mathbb{R}$), the underlying
    Euclidean form $Q_{\mathrm{pos}}(v) = v_0^2 + v_1^2 + v_2^2 + v_3^2$ has zero
    nonzero isotropic vectors.
\end{enumerate}
\end{theorem}
\begin{proof}
Moving each $e_i$ across the preceding factors in $(e_0 e_1 e_2 e_3)^2$ requires
$3 + 2 + 1 = 6$ anticommutations, so $\omega_4^2 = (-1)^6 e_0^2 e_1^2 e_2^2 e_3^2
= \epsilon^4 = 1$ since $\epsilon^2 = 1$. At grade $k = 4$,
$(-1)^4 = (-1)^6 = (-1)^{10} = +1$, so $\alpha(\omega_4) = \tilde{\omega}_4 =
\bar{\omega}_4 = \omega_4$. Hence $\iota(x) = x$ and $x\,\iota(x) = x^2 = 2(1 + \omega_4)$,
whose scalar part is $2 \neq 0$.
\end{proof}

\begin{remark}[Regular-representation determinant $\Nreg(x) = \det(L_x)$]
If one abandons Clifford involution norms altogether and redefines the norm of
$x \in \Cl(V, Q)$ to be the generic finite-dimensional algebra determinant
$\Nreg(x) = \det(L_x)$ (or reduced norm in a semisimple representation), then
$\Nreg(x) = 0 \iff x \text{ is a zero divisor}$ ($x \neq 0$) becomes a tautology valid
in \emph{every} finite-dimensional associative algebra (unrelated to Clifford
structure or conversion constants), while Clauses~B and~C remain false by
Theorem~\ref{thm:disproof-B} and Section~\ref{sec:clause-c}.
\end{remark}

\section{Disproof of Clause C: Minimal Left Ideals vs.\ Spinor Modules}
\label{sec:clause-c}

Clause~C asserts that ``the minimal left ideal of zero divisors is the spinor
module.'' While true for primitive-idempotent left ideals in central simple split
algebras such as $\Cl_{1,1}(F) \cong M_2(F)$ (Theorem~\ref{thm:cl11}), this
identification fails in three distinct structural classes of Clifford algebras.

\begin{theorem}[Degenerate exterior algebras $\Lambda(F^n) = \Cl(F^n, 0)$]
\label{thm:exterior}
Let $F$ be a field with $\mathrm{char}(F) \neq 2$, and consider the totally
degenerate quadratic space $(F^n, 0)$ ($n \ge 1$), whose Clifford algebra is the
exterior algebra $\Lambda(F^n) = \Cl(F^n, [0, \dots, 0])$. Let
$\omega_n = e_0 \wedge \dots \wedge e_{n-1}$ and $M = \Lambda(F^n)\omega_n = F\cdot\omega_n$.
Then:
\begin{enumerate}
  \item $M$ consists of zero divisors (together with $0$) and is the \textbf{unique}
    minimal left ideal of $\Lambda(F^n)$: every nonzero left ideal $J \subseteq \Lambda(F^n)$
    contains $M$;
  \item every generating vector $v \in F^n$ acts identically by \textbf{zero} on $M$
    ($v \wedge \omega_n = 0$), so $M$ is unfaithful ($\ker\rho_M = \Lambda^+(F^n)$)
    and $\dim_F M = 1 < 2^{\lfloor n/2\rfloor}$ for all $n \ge 2$;
  \item $\Lambda(F^n)$ is local with Jacobson radical $\Lambda^+(F^n)$ and possesses
    \textbf{no nontrivial idempotents}: $e^2 = e \implies e \in \{0, 1\}$. Thus $M$
    is not generated by any primitive idempotent and is not a direct summand of
    $\Lambda(F^n)$.
\end{enumerate}
\end{theorem}
\begin{proof}
Since $e_i \wedge \omega_n = 0$ for all $i$, $\Lambda^+(F^n)\omega_n = 0$, so
$\Lambda(F^n)\omega_n = F\cdot\omega_n$ and $\omega_n^2 = 0$. Given any nonzero
$z = \sum_S z_S e_S \in J$, pick a minimal-cardinality subset $S$ with $z_S \neq 0$;
left-multiplying $z$ by the complementary blade $e_{\{0,\dots,n-1\}\setminus S}$
annihilates every $e_T$ with $T \cap (\{0,\dots,n-1\}\setminus S) \neq \emptyset$
(hence every $T \neq S$ with $|T| \ge |S|$), producing $\pm z_S \omega_n \in J$,
whence $\omega_n \in J$ and $M \subseteq J$. Finally, if $e = p + m$ ($p \in F$,
$m \in \Lambda^+(F^n)$) satisfies $e^2 = e$, taking grade $0$ gives $p^2 = p$, so
$p \in \{0, 1\}$. Since $m^{n+1} = 0$: if $p = 0$, $m = m^2 = \dots = m^{n+1} = 0$;
if $p = 1$, $1 + 2m + m^2 = 1 + m \implies m(1 + m) = 0$, and multiplying by the
unit $\sum_{k=0}^n (-m)^k = (1 + m)^{-1}$ gives $m = 0$.
\end{proof}

\begin{theorem}[Split semisimple Clifford algebras and non-unique/unfaithful ideals]
\label{thm:split}
Let $F$ be a field with $\mathrm{char}(F) \neq 2$, and consider
$\Cl_{1,0}(F) = \Cl(F, [1]) \cong F \oplus F$ with central orthogonal idempotents
$p_+ = \frac{1}{2}(1 + e_0)$ and $p_- = \frac{1}{2}(1 - e_0)$. Then:
\begin{enumerate}
  \item $I_+ = \Cl_{1,0}(F)p_+ = F p_+$ and $I_- = \Cl_{1,0}(F)p_- = F p_-$ are
    both minimal left ideals whose nonzero elements are zero divisors;
  \item neither $I_+$ nor $I_-$ is a faithful $\Cl_{1,0}(F)$-module
    ($p_- I_+ = 0$ and $p_+ I_- = 0$);
  \item $I_+$ and $I_-$ are \textbf{not isomorphic} as left $\Cl_{1,0}(F)$-modules,
    since $e_0$ acts as $+1$ on $I_+$ and as $-1$ on $I_-$ (they are distinct
    half-pinor representations, whereas spinor modules of $\Cl^0_{1,0}(F) = F$
    contain no zero divisors at all);
  \item $p_+$ and $p_-$ are zero divisors whose sum $p_+ + p_- = 1$ is a unit, so
    the set of zero divisors (union $\{0\}$) does not form a left ideal.
\end{enumerate}
Analogously, in $\Cl_{0,3}(F)$ ($\cong \mathbb{H} \oplus \mathbb{H}$ for $F = \mathbb{R}$),
$q_\pm = \frac{1}{2}(1 \pm e_0 e_1 e_2)$ generate two non-isomorphic, unfaithful
$4$-dimensional minimal left ideals distinguished by $\omega_3 \mapsto \pm 1$.
\end{theorem}
\begin{proof}
Since $\dim_F I_\pm = 1$, both are minimal; $p_+ p_- = \frac{1}{4}(1 - e_0^2) = 0$
gives zero-divisor membership and unfaithfulness. If $\varphi : I_+ \to I_-$ were
a $\Cl_{1,0}(F)$-linear isomorphism, then for $y = \varphi(p_+) \in I_-$ we would have
$y = \varphi(e_0 p_+) = e_0 \varphi(p_+) = e_0 y = -y$, so $2y = 0 \implies y = 0$,
contradicting bijectivity since $p_+ \neq 0$.
\end{proof}

\begin{corollary}[Division Clifford algebras have no zero-divisor left ideals]
\label{cor:division}
Over any ordered field $F$ ($\mathbb{Q}, \mathbb{R}$), $\Cl_{0,1}(F)$ and
$\Cl_{0,2}(F)$ possess no nonzero left ideals consisting of zero divisors, even
though $\Cl_{0,1}(\mathbb{R}) \cong \mathbb{C}$ and
$\Cl_{0,2}(\mathbb{R}) \cong \mathbb{H}$ possess well-defined spinor modules
($\mathbb{C}$ and $\mathbb{H}$ themselves).
\end{corollary}

\section{Synthesis Across Conventions and Formal Verification}

Table~\ref{tab:summary} synthesizes the status of Clauses~A, B, and~C across
representative Clifford algebras in dimensions $1 \le n \le 4$.

\begin{table}[ht]
\centering
\footnotesize
\setlength{\tabcolsep}{4.2pt}
\renewcommand{\arraystretch}{1.22}
\begin{tabular}{@{}llllll@{}}
\toprule
Algebra ($\mathbb{R}$) & Signature $\sigma$ & Isotropic $v$? & Zero div.? & $x\bar{x}=0$ crit.? & Minimal left ideals \\
\midrule
$\Cl_{0,1} \cong \mathbb{C}$ & $[-1]$ & No & No & Holds ($n=1$) & None of ZDs (div.) \\
$\Cl_{1,0} \cong \mathbb{R}\oplus\mathbb{R}$ & $[+1]$ & \textbf{No} & \textbf{Yes} ($1\pm e_0$) & Holds ($n=1$) & \textbf{2 non-iso.\ unfaithful} \\
$\Cl_{0,2} \cong \mathbb{H}$ & $[-1,-1]$ & No & No & Holds ($n=2$) & None of ZDs (div.) \\
$\Cl_{2,0} \cong M_2(\mathbb{R})$ & $[+1,+1]$ & \textbf{No} & \textbf{Yes} ($1\pm e_0$) & Holds ($n=2$) & Faithful spinor ($\mathbb{R}^2$) \\
$\Cl_{1,1} \cong M_2(\mathbb{R})$ & $[+1,-1]$ & Yes ($e_0\pm e_1$) & Yes ($\ker\cap V=0$) & Holds ($n=2$) & Faithful spinor ($\mathbb{R}^2$) \\
$\Lambda(\mathbb{R}^2)$ & $[0,0]$ & Yes ($V\setminus 0$) & Yes (nilp.) & Holds ($n=2$) & \textbf{1D trivial $V$-act.} \\
$\Cl_{0,3} \cong \mathbb{H}\oplus\mathbb{H}$ & $[-1,-1,-1]$ & \textbf{No} & \textbf{Yes} ($1\pm\omega_3$) & \textbf{Fails} ($x\bar{x}=2x$) & \textbf{2 non-iso.\ unfaithful} \\
$\Cl_{4,0} \cong M_2(\mathbb{H})$ & $[+1,+1,+1,+1]$ & \textbf{No} & \textbf{Yes} ($1\pm\omega_4$) & \textbf{Fails all $\iota$} & Faithful spinor ($\mathbb{H}^2$) \\
$\Cl_{0,4} \cong M_2(\mathbb{H})$ & $[-1,-1,-1,-1]$ & \textbf{No} & \textbf{Yes} ($1\pm\omega_4$) & \textbf{Fails all $\iota$} & Faithful spinor ($\mathbb{H}^2$) \\
\bottomrule
\end{tabular}
\caption{Summary of zero-divisor, norm, isotropy, and ideal properties across
standard Clifford algebras ($e_i^2 = \sigma_i$). Bold entries mark direct
counterexamples to Clauses~A, B, or~C.}
\label{tab:summary}
\end{table}

\subsection*{Computational and Formal Verification}

\begin{itemize}
  \item \textbf{Exact Rational Python Suite (\texttt{verify.py}):} Uses Python~3's
    \texttt{fractions.Fraction} (zero floating-point arithmetic) to verify basis
    isomorphism between bitmask and recursive Clifford models across $16$
    signatures, check exact determinants $\det(L_x)$ and nullspaces on $>1{,}200$
    vectors and multivectors, and confirm every counterexample in
    Theorems~\ref{thm:disproof-B}--\ref{thm:split}.
  \item \textbf{Core Lean~4.19.0 Formalization (\texttt{lean/}):} Compiled with
    flag \texttt{-DwarningAsError=true}, zero external dependencies, and zero
    \texttt{sorry}, \texttt{admit}, \texttt{native\_decide}, or custom axioms.
    Module \texttt{Basic.lean} axiomatizes \texttt{CommRingOps}, \texttt{FieldOps},
    \texttt{CharNeTwoField}, and \texttt{OrderedField} (covering both $\mathbb{Q}$
    and $\mathbb{R}$) and constructs \texttt{QRat} ($\mathbb{Q}$) explicitly as a
    quotient of $\mathbb{Z} \times \mathbb{Z}_{>0}$. Module \texttt{Clifford.lean}
    formalizes Proposition~\ref{prop:recursive} in arbitrary dimension $n$. Module
    \texttt{PositiveRegimes.lean} proves Theorems~\ref{thm:grade1}--\ref{thm:cl11}.
    Finally, \texttt{Disproof.lean} proves
    Theorems~\ref{thm:disproof-B}--\ref{thm:split}, culminating in the master
    theorem
    \begin{center}
      \texttt{conjecture\_00000007622\_false\_over\_ordered\_field}
    \end{center}
    and its unconditional ground specialization
    \begin{center}
      \texttt{conjecture\_00000007622\_disproved\_over\_QRat}.
    \end{center}
\end{itemize}

\begin{thebibliography}{9}
\bibitem{Chevalley1954}
C.~Chevalley, \emph{The Algebraic Theory of Spinors}, Columbia University Press,
New York, 1954.

\bibitem{LawsonMichelsohn1989}
H.~B.~Lawson and M.-L.~Michelsohn, \emph{Spin Geometry}, Princeton Mathematical
Series, vol.~38, Princeton University Press, Princeton, NJ, 1989.

\bibitem{Lounesto2001}
P.~Lounesto, \emph{Clifford Algebras and Spinors}, 2nd ed., London Mathematical
Society Lecture Note Series, vol.~286, Cambridge University Press, Cambridge, 2001.

\bibitem{ABS1964}
M.~F.~Atiyah, R.~Bott, and A.~Shapiro, \emph{Clifford modules}, Topology
\textbf{3} (1964), suppl.~1, 3--38.
\end{thebibliography}

\end{document}
